% sections/sec7_physics.tex
\section{物理学中的张量}
\label{sec:physics}

在物理学与工程学中，张量通常不是通过多重线性映射来定义，而是通过它们在不同坐标系下的变换行为来定义的.以下是在许多物理学教科书中常见的一种张量定义：

设 $V$ 为 $n$ 维向量空间（$n \geq 1$）.在 $V$ 上，$0$ 阶张量就是一个标量.对于 $k \geq 1$，$V$ 上的一个 \newterm{contravariant-tensor}{逆变张量}{contravariant tensor}（$k$ 阶）是一个在 $V$ 的每个坐标系下均具有 $n^k$ 个分量的对象 $T$；如果 $\{T^{i_1, \dots, i_k}\}_{1 \leq i_1, \dots, i_k \leq n}$ 与 $\{\widetilde{T}^{i_1, \dots, i_k}\}_{1 \leq i_1, \dots, i_k \leq n}$ 是 $T$ 在 $V$ 的两个坐标系下的分量，则它们满足变换律：
\begin{equation}
  \widetilde{T}^{i_1, \dots, i_k} = \sum_{1 \leq j_1, \dots, j_k \leq n} T^{j_1, \dots, j_k} a_{i_1 j_1} \dots a_{i_k j_k}, \label{eq:7.1}
\end{equation}
其中 $(a_{ij})$ 是将 $V$ 的第一个坐标系用第二个坐标系表示的变换矩阵.简言之，逆变 $k$ 阶张量就是``按规则~\eqref{eq:7.1} 变换的量''.

这里用分量所描述的对象，本质上就是张量幂 $V^{\otimes k}$ 中的一个元素.事实上，所谓``坐标系''就是 $V$ 的一组基.设 $\{e_1, \dots, e_n\}$ 是 $V$ 的一组基，在其中 $T$ 具有分量 $\{T^{i_1, \dots, i_k}\}$，将这些分量作为 $V^{\otimes k}$ 的基 $\{e_{i_1} \otimes \dots \otimes e_{i_k}\}$ 的展开系数：
\[
  \sum_{1 \leq i_1, \dots, i_k \leq n} T^{i_1, \dots, i_k} e_{i_1} \otimes \dots \otimes e_{i_k}.
\]
若将其用另一组基（``新坐标系''）$\{f_1, \dots, f_n\}$ 展开，设 $e_j = \sum_{i=1}^n a_{ij} f_i$，代入并重排求和项：
\begin{align*}
  &\sum_{1 \leq j_1, \dots, j_k \leq n} T^{j_1, \dots, j_k} e_{j_1} \otimes \dots \otimes e_{j_k} \\
  &= \sum_{1 \leq j_1, \dots, j_k \leq n} T^{j_1, \dots, j_k} \left(\sum_{i_1=1}^n a_{i_1 j_1} f_{i_1}\right) \otimes \dots \otimes \left(\sum_{i_k=1}^n a_{i_k j_k} f_{i_k}\right) \\
  &= \sum_{1 \leq i_1, \dots, i_k \leq n} \left(\sum_{1 \leq j_1, \dots, j_k \leq n} T^{j_1, \dots, j_k} a_{i_1 j_1} \dots a_{i_k j_k}\right) f_{i_1} \otimes \dots \otimes f_{i_k} \\
  &= \sum_{1 \leq i_1, \dots, i_k \leq n} \widetilde{T}^{i_1, \dots, i_k} f_{i_1} \otimes \dots \otimes f_{i_k}.
\end{align*}
因此在物理学中，$V$ 上逆变 $k$ 阶张量的分量，正是 $V^{\otimes k}$ 中的元素在 $V^{\otimes k}$ 的某组基下的展开坐标.在物理学中，$V$ 的维数通常为 $3$ 或 $4$.

现在从 $V$ 的张量幂转向其对偶空间 $V^\vee$ 的张量幂.我们要比较 $(V^\vee)^{\otimes \ell}$ 中的元素在由对偶基构成的坐标系下的表示.

为了与物理学和微分几何中的习惯一致，从现在起将基 $\{e_1, \dots, e_n\}$ 与 $\{f_1, \dots, f_n\}$ 的对偶基记作 $\{e^1, \dots, e^n\}$ 与 $\{f^1, \dots, f^n\}$（指标写在上标），满足 $e^i(e_j) = f^i(f_j) = \delta_{ij}$.$V$ 的一般元素写为 $\sum c^i e_i$（基带指标下角标，系数带指标上角标），而对偶空间 $V^\vee$ 的一般元素写为 $\sum c_i e^i$（基带指标上角标，系数带指标下角标）.

任取张量 $\mathsf{T} \in (V^\vee)^{\otimes \ell}$，在基 $\{e^{i_1} \otimes \dots \otimes e^{i_\ell}\}$ 下展开为
\begin{equation}
  \mathsf{T} = \sum_{1 \leq i_1, \dots, i_\ell \leq n} T_{i_1, \dots, i_\ell} e^{i_1} \otimes \dots \otimes e^{i_\ell}. \label{eq:7.2}
\end{equation}
已知在 $V$ 中 $e_j = \sum_{i=1}^n a_{ij} f_i$.计算对偶基的变换关系：
\begin{equation}
  e_j = \sum_{i=1}^n a_{ij} f_i \quad\Longrightarrow\quad f^j = \sum_{i=1}^n a_{ji} e^i. \label{eq:7.3}
\end{equation}
注意式~\eqref{eq:7.3} 中出现了转置矩阵 $(a_{ji})$.将矩阵 $(a_{ij})$ 的逆矩阵记作 $(a^{ij})$，则对偶基的逆向变换为：
\begin{equation}
  e^j = \sum_{i=1}^n a^{ji} f^i. \label{eq:7.6}
\end{equation}
将式~\eqref{eq:7.6} 代入式~\eqref{eq:7.2}，即得分量的协变变换律：
\begin{equation}
  \widetilde{T}_{i_1, \dots, i_\ell} = \sum_{1 \leq j_1, \dots, j_\ell \leq n} T_{j_1, \dots, j_\ell} a^{j_1 i_1} \dots a^{j_\ell i_\ell}. \label{eq:7.7}
\end{equation}
按规则~\eqref{eq:7.7} 变换的量称为 $\ell$ 阶\newterm{covariant-tensor}{协变张量}{covariant tensor}.它无非就是对偶张量空间 $(V^\vee)^{\otimes \ell}$ 中的元素.

更一般地，按如下规则变换的多重指标量：
\begin{equation}
  \widetilde{T}_{j_1 \dots j_\ell}^{i_1 \dots i_k} = \sum_{\substack{1 \leq p_1, \dots, p_k \leq n \\ 1 \leq q_1, \dots, q_\ell \leq n}} T_{q_1 \dots q_\ell}^{p_1 \dots p_k} a_{i_1 p_1} \dots a_{i_k p_k} a^{q_1 j_1} \dots a^{q_\ell j_\ell} \label{eq:7.8}
\end{equation}
称为\newterm{type-k-l-tensor}{$(k, \ell)$ 型张量}{tensor of type (k, l)}（其总阶数为 $k + \ell$）.它正是混合张量空间
\[
  V^{\otimes k} \otimes (V^\vee)^{\otimes \ell}
\]
中的元素.

\begin{example}\label{ex:rank2-tensor-physics}
$2$ 阶张量分为三类：
\begin{enumerate}
  \item \textbf{$(2, 0)$ 型张量}：属于 $V^{\otimes 2}$，其矩阵变换规则为 $(\widetilde{T}^{ij}) = A (T^{ij}) A^\top$；
  \item \textbf{$(0, 2)$ 型张量}：属于 $(V^\vee)^{\otimes 2}$，其矩阵变换规则为 $(\widetilde{T}_{ij}) = (A^{-1})^\top (T_{ij}) A^{-1}$，它对应于 $V$ 上的\textbf{双线性型}（例如内积度规、闵氏度规）；
  \item \textbf{$(1, 1)$ 型张量}：属于 $V \otimes V^\vee \cong \End(V)$，其矩阵变换规则为 $(\widetilde{T}^i_j) = A (T^i_j) A^{-1}$，它对应于 $V$ 上的\textbf{线性变换}（相似变换）.
\end{enumerate}
\end{example}

\begin{remark}\label{rem:einstein-quote}
数学家与物理学家看待张量的角度截然不同：
\begin{itemize}
  \item \textbf{数学家}：张量是张量空间（由多重线性映射普遍性质定义的模或向量空间）中的元素，完全独立于任何基或坐标系.
  \item \textbf{物理学家}：``张量是由一个或多个指标组织的分量体系，在坐标变换下服从特定的变换规则''\cite{feynman-lectures}.
\end{itemize}

Einstein 在其 1916 年关于广义相对论的基础论文中写道 \cite[第~157~页]{einstein-gr}：
\begin{quote}
``设对于任意坐标系，某些对象（'张量'）由若干坐标函数定义，这些函数称为张量的'分量'.当从原坐标系变换到新坐标系且变换已知时，存在确定的法则来计算新坐标系下的这些分量.被称为张量的对象其特征在于其分量的变换方程是线性且齐次的.因此，若所有分量在原坐标系中全为零，则它们在新坐标系中亦全部为零.''
\end{quote}
Einstein 这里所说的``线性且齐次变换''，在现代数学语言中就是多重线性变换.

Hassler Whitney 在回忆早年学习张量的困惑时曾生动地写道 \cite[第~114~页]{whitney-history}：
\begin{quote}
``我不得不摆弄张量；但我连见都不被允许见到它们，只能被告知它们在坐标变换下不断变换的外衣装扮，我又怎能掌握它们呢？我无论如何都必须一把揪出这些淘气的小家伙，直面它们的真面目！''
\end{quote}
\end{remark}

\begin{example}\label{ex:stress-tensor}
在连续介质力学中，最基础的 $2$ 阶张量是柯西应力张量.在一点处施加力时，产生的应力方向未必沿着力的方向，且这种响应在输入力上是线性的.因此某点处的应力由一个线性变换描述，从而是一个 $(1, 1)$ 型张量（属于 $V \otimes V^\vee$）.由于受力情况在物体各点处不同，应力实际上是由各点处张量构成的光滑族，即\newterm{tensor-field}{张量场}{tensor field}.

在微分流形 $M$ 上，张量场可以在局部定义（在每点 $p \in M$ 选取切空间与余切空间张量积 $T_p(M)^{\otimes k} \otimes T_p(M)^{\vee \otimes \ell}$ 中的元素并随 $p$ 光滑变化），亦可在全局定义为光滑切向量场模 $\mathfrak{X}(M)$ 在光滑函数环 $C^\infty(M)$ 上的张量积：
\[
  \mathfrak{X}(M)^{\otimes k} \otimes_{C^\infty(M)} \mathfrak{X}(M)^{\vee \otimes \ell}.
\]
此时局部坐标变换 $x \mapsto y$ 诱导的变换矩阵正是雅可比偏导数矩阵 $\frac{\partial y^i}{\partial x^j}$.
\end{example}

需要注意，\textbf{并非所有带指标的量都是张量}.例如微分几何中的克里斯托费尔符号（Christoffel symbols）$\Gamma^k_{ij}$，虽然带有三个指标，但在坐标变换下包含偏导数的二次非齐次项，不满足多重线性变换法则，因此它不是张量.

在量子力学中，张量积扮演着更为根本的角色.在经典力学中，复合力学系统的相空间是子系统流形的笛卡尔积 $M \times N$；而在量子力学中，系统的状态由复 Hilbert 空间（如 $L^2(\R^3)$）中的单位射线描述.当组合两个量子系统时，复合系统的状态空间是各自 Hilbert 空间的张量积：
\[
  \mathcal{H}_1 \otimes_\C \mathcal{H}_2.
\]
在直积 $M \times N$ 中，每个点都是形如 $(x, y)$ 的平凡对；但在张量积 $\mathcal{H}_1 \otimes_\C \mathcal{H}_2$ 中，绝大多数张量都\textbf{不是简单张量}.那些无法表示为单一简单张量 $v \otimes w$ 的态，正是量子力学中最神奇也最深刻的现象——\newterm{entanglement}{量子纠缠态}{entangled states}.
